\begin{scanpage}{0227}{209}
If the receiving antenna operates with the same polarization as the transmitting antenna, then the corresponding value of the effective area will equal
\[\sigma_a(\vartheta,\vartheta_0)=\frac{\lambda^2}\pi\cos^4\alpha\cdot|F(\vartheta,\vartheta_0)|^2\tag{34.11}\]
In the case of radar when the transmitting and receiving antennas are combined and the polarization is arbitrary, the vibrator's scattering properties are characterized by the average value
\[\bar\sigma(\vartheta)=\frac1{2\pi}\int_0^{2\pi}\sigma_a(\vartheta,\vartheta_0)d\alpha=\frac{3\lambda^2}{8\pi}|F(\vartheta,\vartheta_0)|^2.\tag{34.12}\]
In Figures 77 and 78, the results of calculations performed on the basis of Equations (34.12) and (32.08) for the case of normal irradiation $\left(\vartheta=\frac\pi2\right)$ are represented by the dotted lines. Figure 77 illustrates the dependence of the function $\bar\sigma$ on the quantity $kL$ with a given value of $\Omega_p=2\ln\frac{2L}a=15$,---that is, when the ratio of the vibrator's length to its diameter equals $L/2a=452$. In Figure 78 the graph of the function $\bar\sigma$ is constructed as a function of the frequency $f=c/\lambda\cdot10^{-6}$ (in megahertz) for the prescribed parameters $L=5$ cm and $\Omega_p=15$. Here the curves plotted by Lindroth [79] are drawn with a continuous line, and the curve in Figure 77 calculated by Van Vleck et al.\ [86] is traced by the dash-dot-curve.

The curves of Lindroth and Van Vleck were calculated by integrating the current which is found as a result of the approximate solution of the integral equation. However, this procedure was performed in [79] and [86] in a different way. Lindroth obtained an expression for the fringing field in the form of an expansion in reciprocal powers of the parameter $\Omega_p$. The expression includes terms of the order of $\Omega_p^{-3}$. In [86] a different kind of approximation was used which led, as can be seen from Figure 77, to rather rough results especially in the resonance region. Our curve (the dotted area) agrees almost everywhere within the limits of graphical precision with the curve of Lindroth. A noticeable divergence is observed only in the magnitude of the first resonance peak.
\end{scanpage}
